\documentclass{fmfkpb}

\avtor{Luka Lavš}
\naslov{Mooreovi grafi in posplošeni večkotniki}
\letnica{2026}
\mentor{doc. dr. Janoš Vidali}

\newcommand{\red}{\operatorname{red}}
\newcommand{\srg}{\operatorname{srg}}
\newcommand{\diam}{diam}
\newcommand{\N}{\mathbb{N}}
\newcommand{\ZZ}{\mathbb{Z}}
\newcommand{\reach}{\operatorname{Reach}}
\newcommand{\gr}{\mathscr{G}}

\newcommand{\A}{\mathcal{A}}
\newcommand{\II}{\mathbf{I}}
\newcommand{\IIw}[2]{#1 \;\II\; #2}
\newcommand{\IIwd}[2]{#1 \;\II^D\; #2}
\newcommand{\PP}{\mathcal{P}}
\newcommand{\LL}{\mathcal{L}}

\begin{document}

\makefrontpage
\cleardoublepage

\section{Mooreovi grafi in posplošeni večkotniki}

\subsection{Uvodne definicije}

    \begin{itemize}
        \item[\textbullet] Graf $G$ ima \textbf{premer} $D$, če je največja razdalja med  
        katerima koli dvema vozliščema enaka $D$.
        \item[\textbullet] Graf $G$ ima \textbf{maksimalno stopnjo} $\Delta$, če je
        $\Delta = \max_{v \in V(G)} \deg(v)$.
        \item[\textbullet] Graf $G$ ima \textbf{ožino} $g$, če je dolžina 
        najkrajšega cikla v $G$ enaka $g$.
        \item[\textbullet] Graf $G$ je \textbf{$\Delta$-regularen}, če ima vsako 
        vozlišče stopnjo $\Delta$. 
    \end{itemize}

Sedaj pa vzemimo graf danega premera $D$ 
ter maksimalne stopnje $\Delta$, ter skonstruirajmo zgornjo mejo za število 
njegovih vozlišč. 

Vzemimo neko vozlišče in ga povežimo z največ $\Delta$ drugimi, te 
vozlišča pa nato povežimo z največ $\Delta - 1$ novimi vozlišči, to 
nadaljujemo dokler naš drevesni podgraf ne doseže premera $D$. 

Tako za zgornjo mejo števila vozlišč dobimo
$1 + \Delta + \Delta(\Delta - 1) + \ldots + \Delta(\Delta - 1)^{D - 1}$.
Kar je seveda enako $\frac{\Delta(\Delta - 1)^D - 2}{\Delta - 2}$. Označimo 
ta rezultat z $M_{D, \Delta}$. 

\subsection{Mooreov graf}
Grafu, ki tako zgornjo mejo doseže pravimo Mooreov graf.  

\subsection{Lastnosti Mooreovih grafov}
Po konstrukciji zgornje meje se s krajšim dokazom prepričamo, da 
so taki grafi $\Delta$-regularni. Namreč zgornjo mejo za 
število vozlišč lahko izpeljemo iz poljubnega vozlišča, če pa 
bi konstrukcijo začeli iz vozlišča stopnje manj kot $\Delta$, 
prave zgornje meje ne bi dosegli. $
    1 + \deg(u) + \deg(u) (\Delta -1) + \ldots \deg(u) (\Delta -1)^{D - 1}
    < M_{D, \Delta}$ vodi v protislovje. 
\newline 
Prav tako se iz konstrukcije izpelje dejstvo, da 
med poljubnima vozliščema obstaja \emph{natanko ena pot
 dolžine največ $D$}.
To je res, ker konstrukcija predpostavi, da je vsako 
vozlišče koren nekega drevesnega podgrafa, kjer so poti enolične 
in dolžine največ $D$.
\newline 
Podobno se izkaže tudi, da je ožina Mooreovega grafa enaka 
$g=2D + 1$.

\subsection{Matrična enačba Mooreovih grafov}
Dejstvo o tem, da med poljubnima vozliščema obstaja natanko ena pot 
dolžine največ $D$, nam omogoča zapis matrične enačbe, kateri 
morajo Mooreovi grafi zadoščati. 

Matrika izpeljana rekurzivno s predpisom 
$G_{\Delta, i} = A G_{\Delta, i - 1} - (\Delta - 1) G_{\Delta, i - 2}$
($2 \le i \le D$),
$G_{\Delta, 1} = A + I$, $G_{\Delta, 0} = I$, kjer je $A$ matrika sosednosti
grafa, namreč meri število poti dolžin največ $i$ med poljubnima
vozliščema. 

Z njeno pomočjo pa pridemo do enačbe $G_{\Delta, D} = J$, 
kjer je $J$ matrika samih enic.
\newline 
Za primer ko je premer grafa enak $D=2$ dobimo enačbo 
$A^2 + A - (\Delta - 1) I = J$. Prek nje pa omejitve za 
lastne vrednosti in posledično tudi 
omejitve za možne maksimalne stopnje Mooreovih grafov.
Tako se izkaže, da so možne stopnje za premer $D=2$ enake 
$\{2, 3, 7, 57\}$. 
\newline 
Mooreovi grafi s stopnjami $2, 3$ in $7$ so petkotnik, Petersenov graf
ter Hoffman-Singeltonov graf. Obstoj grafa stopnje $57$ 
pa je še vedno odprto vprašanje.

\subsection{Redkost ter razširitve}
Zahtevke o omejenem premeru, maksimalni stopnji, regularnosti, 
ožini\\ so kar "močne". Zato so Mooreovi grafi zelo redki. Preprosti 
poizkusi konstrukcij s pohlepnimi algoritmi seveda običajno spodletijo. 
Tako je na primer mogoče Petersenov graf zgraditi s takšnim
algoritmom, medtem ko Hoffman-Singletonov graf zahteva bolj
sofisticirano konstrukcijo. 

Nekateri avtorji so zato predlagali analize "blizu" ~Mooreovih grafov, 
kjer so zahteve nekoliko omiljene. Tej iščejo grafe, ki dosegajo 
Mooreovo mejo z majhnim defektom $\left(M_{\Delta, D} - \delta\right)$.

Nekateri pa so se lotili tudi problema z dano ožino $g$ ter 
maksimalno stopnjo $D$. Tu je zanimivo, da se spodnja meja števila 
vozlišč za tak problem ujema z zgornjo mejo Mooreovih grafov (
    v primeru ožine oblike $g = 2D + 1$
). Zato bomo rekli, da tak problem razširja Mooreove grafe. 

Če pa se v razširitvenem problemu omejimo na sode ožine, lahko 
opazimo, da so taki grafi natanko incidenčni grafi 
posplošenih večkotnikov.

\subsection{Geometrija in posplošeni večkotniki}
Na kratko, geometrija ranga 2 je trojica$\left(\PP, \LL, \II\right)$,
kjer sta $\PP$ in $\LL$ neprazni disjunktni množici, ki predstavljata 
točke in premice, medtem ko je $\II$ relacija incidenčnosti nad njima.

Posplošeni n-kotnik pa je geometrija ranga 2, ki kot podgeometrijo 
ne vsebuje nobenega večkotnika dolžine manjše od $n$; ki poljubna 
dva elementa unije točk in premic povezuje z nekim običajnem $n$-kotnikom 
(kot podgeometrijo); ter ki ima vsaj en $(n+1)$-kotnik kot podgeometrijo.

Incidenčni graf posplošenega večkotnika pa je nato le graf, z unijo 
točk in premic kot množico vozlišč, ter povezavami določenimi prek 
incidenčne relacije.

\end{document}